\color{unimain}{\subsection{Защита от внутренних повреждений батареи статических конденсаторов}}
\color{black}

\begin{enumerate}

\item
Защита от внутренних повреждений (небалансная защита) предназначена для защиты батареи статических конденсаторов при замыкании ряда или пробое единичных конденсаторов. Защита контролирует токи небаланса, которые протекают в цепях проводников, соединяющих средние точки параллельных ветвей каждой фазы БСК. В случае равномерного повреждения конденсаторов параллельных ветвей в любой из фаз БСК возможен отказ данной защиты. 

\item 
Уставки ЗВП приведены в таблице \ref{app_set:zvp}. Алгоритм работы ЗВП приведен на рисунке \ref{pict:ZVP}.

\medskip
\begin{figure}[H]
\centering
\includegraphics[width=1\textwidth,height=1\textheight,keepaspectratio]{\fbpath/09.01 БСК/261.0231 ЗВП/img/ZVP.pdf}
\captionof{figure}{Функциональная схема логики ступеней ЗВП}\label{pict:ZVP}
\end{figure}

\item
Естественный ток небаланса, протекающий в цепи проводника, соединяющего средние точки параллельных ветвей фазы БСК, может быть вызван: 
\begin{itemize}
\item разбросом параметров конденсаторов в пределах допустимого значения, регламентируемого заводом-изготовителем; 
\item изменением емкости под воздействием температуры окружающей среды (ввиду изменения свойств диэлектрика); 
\item естественным старением конденсаторов. 
\end{itemize}

\item
Алгоритмом предусматривается предварительная компенсация естественного тока небаланса в фазах, выполненная в соответствии с формулами: 
\begin{equation}
\label{eq:INB}
I_{\textsl{НБ.х}} = \mid I_{\textsl{СТ.х}} - K_{\textsl{комп.х}} \cdot I_{\textsl{ВВ.х}} \mid
\end{equation}

\begin{eqexpl}[10mm]
\item{${\textsl{x}}$} наименование контролируемой фазы (A, B, C); 
\item{$I_{\textsl{НБ.х}}$} значение тока небаланса с учетом компенсации; 
\item{$I_{\textsl{СТ.х}}$} значение тока небаланса, измеренное в цепи проводника, соединяющего средние точки 
параллельных ветвей фазы БСК; 
\item{$I_{\textsl{ВВ.х}}$} значение фазного тока, измеренное со стороны высоковольтных вводов БСК; 
\item{$K_{\textsl{комп.х}}$} коэффициент компенсации тока небаланса (индивидуальный для каждой фазы). 
\end{eqexpl}

\item
Защита содержит сигнальную и отключающие ступени. Каждая из ступеней имеет регулируемое значение по току срабатывания, задаваемое параметром \textbf{«Iсигн ЗВП»} и \textbf{«Iоткл ЗВП»} для сигнальной и отключающей ступени соответственно. Сигнальная ступень с задержкой по времени \textbf{«Тсигн ЗВП»} действует на сигнализацию, отключающая ступень с задержкой по времени \textbf{«Тоткл ЗВП»} действует на  отключение выключателя БСК. 


\end{enumerate} 